% shared_data_catalog.tex -- the "Data and event catalog" Methods text, written once for both papers (LeafletFA and SpliceVI).
% Every number is a ledger macro (paper/ledger/results_ledger.tsv; SpliceVI copies: paper/ledger/for_splicevi_catalog_rows.tsv).
% Sources: CATALOG_PROVENANCE.md, DEFINITIONS.md. No contractions, no em dashes, American spelling. Draft 2026-11-01.

\paragraph{Data.}
We used two full-length single-cell atlases of the mouse: Tabula Muris Senis (TMS), FACS-sorted single cells from 23
tissues of 18 mice aged 3, 18 and 24 months profiled with Smart-seq2 \cite{Tabula_Muris_Consortium2020-ga}, and the Allen
Brain atlas of the mouse isocortex and hippocampus profiled with Smart-seq v4 (GEO GSE185862) \cite{mouse-allen-brain-ssv4}.
For TMS we used the consortium's alignments to GRCm38 with the GENCODE vM19 annotation; Allen reads were aligned to the
same genome and annotation with STAR 2.7.10b in two-pass mode. Cell annotations were the consortium's tissue and cell
ontology class for TMS and the cluster labels of the GEO metadata for Allen. Gene expression was taken from each atlas's
published count matrix (exonic counts for Allen), restricted to protein-coding and long noncoding RNA genes with more than
one exon that are shared by both atlases and carry catalog junctions (\CatalogGenesMouse{} genes), and normalized for mean
transcript length.

\paragraph{Event catalog.}
Splice junction reads were extracted from each cell's alignments with regtools \cite{regtools} (anchor of at least 6 nt,
introns of 50 to 500{,}000 bp, strand from the XS tag) and pooled over all \FlowCellsReadable{} cells of both atlases.
Junctions with at least 100 reads in at least 10 cells and a canonical splice motif (GT-AG, GC-AG or AT-AC) were kept,
and each was assigned the gene of the transcripts whose exon boundaries it matches within 1 bp at either end; junctions
without a gene, or with more than one, were removed. An alternative transcript structure event (ATSE) is a connected
component of a gene's junctions in which two junctions are linked if they share a donor or an acceptor site; a junction
was removed from its event if it carried less than 1\% of the reads at either of its splice sites, and an event needs at
least two junctions, so junctions that share no splice site with another junction (singletons) are not part of the
catalog. A mutually exclusive exon pair forms one ATSE only if another junction links its two halves; otherwise it
appears as two ATSEs, an alternative-3$'$ and an alternative-5$'$ pair. After matching cells to expression and quality
control (\FlowCellsQC{} cells), junctions observed in fewer than 75 cells were removed, ATSEs were re-derived on the
remaining junctions (which left \FlowSingletonsRegroup{} junctions without a partner), and cells with fewer than 500
detected junctions were removed. The catalog contains \CatalogJunctionsMouse{} junctions in \CatalogATSEsMouse{} ATSEs
across \CatalogGenesMouse{} genes, measured in \CatalogCellsMouse{} cells (\CatalogCellsTMS{} from TMS and
\CatalogCellsAllen{} from Allen; Supplementary Table~\ref{S-tab:catalogflow}). A further \FlowAllenOmitted{} Allen cells,
stored in a different folder layout, were missed when the input files were listed and are not in the catalog.

\paragraph{Training and test cells.}
Each atlas was split at random into 70\% training and 30\% test cells, stratified by broad cell type and age (random seed
42); the split is per cell, so the same mice contribute to both sets. Tabula Muris Senis has \SplitTrainTMS{} training and
\SplitTestTMS{} test cells, and Allen \SplitTrainAllen{} and \SplitTestAllen{}.
